\documentclass[12pt]{article}
\usepackage[utf8]{inputenc}
\usepackage[T1]{fontenc}
\usepackage{hyperref}
\title{Kei Language Resources for K-12 Curriculum Development}
\author{Compiled by Claude Sonnet 5.5 (claude-sonnet-5-5) for Andr\'as Kornai}
\date{}
\begin{document}
\maketitle

\section{Introduction}
Kei (Keiese, Veveu Evav; ISO 639-3 kei) is an Austronesian language of the Kei Islands, Maluku, Indonesia. Search results give about 85,000 speakers in one source. Dialects named include Kei Kecil, Kei Besar, Ta'am, Tanimbar Kei and Tayando. The search turned up no dedicated dictionary, grammar or school materials. It also reports no official spelling system; the Latin script was first used by Dutch missionaries.

\section{Primary Educational Resources}
None found.

\section{Linguistic and Reference Materials}
\begin{itemize}
\item Overview: \url{https://en.wikipedia.org/wiki/Kei_language}
\item Omniglot page on Kei (Veveu Evav), with an alphabet chart: \url{https://www.omniglot.com/writing/kei.htm}
\end{itemize}
The Indonesian Wikipedia page on Tanimbar Kei cites a 1990 SIL study, ``Bahasa Kei (Evav): Hasil Penelitian Selama Satu Abad''. The study itself was not seen, so it is not listed as a resource.

\section{Religious and Community Resources}
\begin{itemize}
\item Scripture Earth lists Kei Bible portions: \url{https://www.scriptureearth.org/00eng.php?iso=kei}
\end{itemize}

\section{Digital and Audio Resources}
\begin{itemize}
\item Global Recordings Network page for Kei: \url{https://globalrecordings.net/en/language/kei}
\end{itemize}

\section{Curriculum Development Context}
Kei is spoken by a large community, but the search found no school materials, standard orthography or lexicon. A curriculum would first need an agreed orthography and a documented lexicon. The Bible portions and the Global Recordings material could serve as initial text and audio.
\end{document}